\chapter{Méthodes des différences finies}\label{ch:qm:finite-difference-methods}

Le lendemain d'un raffinement de grille, le gamma d'un nouveau pricer près du strike présente des dents de scie. Les prix ont convergé au centime près~; leur dérivée seconde oscille d'un nœud à l'autre,
positive, négative, positive, avec des écarts vingt fois supérieurs au gamma réel, et les limites du hedger s'allument en rouge sans raison de marché. Raffiner de nouveau la grille double l'amplitude des dents de scie. Rien
ne cloche dans l'équation ni dans le code au sens ordinaire~: le schéma en temps fait exactement ce que son facteur d'amplification annonce. Ce chapitre discrétise les équations rétrogrades
du chapitre~4 sur une grille, analyse la consistance, la stabilité et les oscillations, et construit le solveur que la société utilise pour les problèmes à un et à deux facteurs.

\section{Discrétisation d'une équation parabolique}

D'après la formule de Feynman--Kac, un prix européen $u(\tau, x)$, en fonction du temps restant avant l'échéance $\tau$ et du logarithme du prix $x = \ln S$, vérifie $u_\tau = \mathcal Lu$ avec $\mathcal Lu = au_{xx} + bu_x - cu$, $a =
\frac12\sigma^2$, $b = r - \frac12\sigma^2$, $c = r$, à partir du payoff en $\tau = 0$.

\begin{definition}[Méthode aux différences finies, méthode des lignes, grille non uniforme]\label{def:qm:finite-difference-methods:fd}
Une \emph{méthode aux différences finies}\index{méthode aux différences finies} remplace les dérivées d'une équation différentielle par des quotients aux différences des valeurs sur une grille de nœuds. La \emph{méthode des
lignes}\index{méthode des lignes} discrétise d'abord l'espace, ce qui transforme l'équation en un système d'équations différentielles ordinaires $\dot U = L_hU$ en temps, puis applique un schéma en temps. Sur une
\emph{grille non uniforme}\index{grille non uniforme} les pas $h_i^- = x_i - x_{i-1}$ et $h_i^+ = x_{i+1} - x_i$ diffèrent~; les formules à trois points
\[u_{xx} \approx \frac{2u_{i-1}}{h^-(h^- + h^+)} - \frac{2u_i}{h^-h^+} + \frac{2u_{i+1}}{h^+(h^- + h^+)},\]
\[u_x \approx \frac{-h^+u_{i-1}}{h^-(h^- + h^+)} + \frac{(h^+ - h^-)u_i}{h^-h^+} + \frac{h^-u_{i+1}}{h^+(h^- + h^+)}\]
sont exactes pour les polynômes du second degré.
\end{definition}

\begin{definition}[Erreur de troncature, schémas consistants et stables]\label{def:qm:finite-difference-methods:consistent}
L'\emph{erreur de troncature}\index{erreur de troncature} d'un schéma est le résidu obtenu lorsqu'on y substitue la solution exacte. Un schéma est un \emph{schéma consistant}\index{schéma consistant} si son
erreur de troncature tend vers zéro avec les pas, et un \emph{schéma stable}\index{schéma stable} si ses solutions restent bornées, uniformément en les pas, sur un intervalle de temps fixé.
\end{definition}

\begin{theorem}[Théorème d'équivalence de Lax]\label{thm:qm:finite-difference-methods:lax}
Pour un problème de valeur initiale linéaire bien posé, un schéma aux différences finies consistant converge si et seulement s'il est stable (Lax et Richtmyer, 1956).
\end{theorem}

Le théorème répartit le travail~: la consistance relève du théorème de Taylor, et la stabilité fait l'objet des deux sections suivantes. L'opérateur de la société (\cref{tut:qm:finite-difference-methods:sawtooth}) construit la
matrice tridiagonale $L_h$ sur une grille quelconque~; chaque pas implicite est une résolution par l'\omterm{def:qm:floating-point-and-numerical-linear-algebra:factor}{algorithme de Thomas (matrices tridiagonales)} du chapitre~25, en $O(n)$ opérations. Une grille étirée par $x_i = x^* + \alpha\sinh(\xi_i)$ avec
des $\xi_i$ uniformes, place les nœuds là où le payoff se courbe~: avec $\alpha = 0.05$ autour du strike, elle divise l'erreur sur le prix du call par environ 2.9 à nombre de nœuds égal (0.0079 contre 0.0225 avec 50
nœuds, 0.00012 contre 0.00035 avec 400).

\section{Schémas explicite, implicite et theta}

\begin{definition}[Schémas explicite, implicite, theta et de Crank--Nicolson]\label{def:qm:finite-difference-methods:theta}
Avec un pas de temps $\Delta\tau$, le \emph{theta-schéma}\index{theta-schéma} s'écrit $(I - \theta\Delta\tau L_h)U^{n+1} = (I + (1 - \theta)\Delta\tau L_h)U^n$. Il est le \emph{schéma explicite}\index{schéma explicite}
(Euler progressif) pour $\theta = 0$, le \emph{schéma implicite}\index{schéma implicite} (Euler rétrograde) pour $\theta = 1$ et le \emph{schéma de Crank--Nicolson}\index{schéma de Crank--Nicolson} (Crank et Nicolson,
1947) pour $\theta = \frac12$~; son erreur en temps est en $O(\Delta\tau)$ sauf pour $\theta = \frac12$, où elle est en $O(\Delta\tau^2)$.
\end{definition}

\begin{definition}[Analyse de stabilité de von Neumann, condition CFL]\label{def:qm:finite-difference-methods:vn}
L'\emph{analyse de stabilité de von Neumann}\index{analyse de stabilité de von Neumann} substitue un mode de Fourier $U_j^n = g^ne^{\mathrm ij\xi}$ dans un schéma à coefficients constants et exige $|g(\xi)| \le 1 +
O(\Delta\tau)$ pour toute fréquence. La \emph{condition CFL}\index{condition CFL} (Courant, Friedrichs et Lewy, 1928) est la restriction qui en résulte sur le pas de temps en fonction du pas d'espace, pour un
\omterm{def:qm:finite-difference-methods:theta}{schéma explicite}.
\end{definition}

\begin{proposition}[Facteurs d'amplification des theta-schémas]\label{prop:qm:finite-difference-methods:amp}
Pour $u_\tau = au_{xx}$ sur une grille uniforme, avec $\lambda = a\Delta\tau/\Delta x^2$ et $s = \sin^2(\xi/2)$,
\[g(\xi) = \frac{1 - 4(1 - \theta)\lambda s}{1 + 4\theta\lambda s}.\]
Le \omterm{def:qm:finite-difference-methods:theta}{schéma explicite} est stable si et seulement si $\lambda \le \frac12$, c'est-à-dire $\Delta\tau \le \Delta x^2/\sigma^2$~; les \omterm{def:qm:finite-difference-methods:theta}{schémas implicite} et de Crank--Nicolson sont stables pour tout $\lambda$.
\end{proposition}

\begin{proof}
$U_{j+1} - 2U_j + U_{j-1} = (e^{\mathrm i\xi} - 2 + e^{-\mathrm i\xi})U_j = -4sU_j$, et la substitution dans le schéma donne $g$. En $\xi = \pi$, le facteur explicite vaut $1 - 4\lambda$, qui est au moins égal à $-1$ si et seulement si
$\lambda \le \frac12$. Pour $\theta \ge \frac12$, le numérateur est au moins égal à $-1$ fois le dénominateur pour tout $\lambda$, donc $|g| \le 1$.
\end{proof}

La limite est nette. Avec 200 nœuds sur la grille, le \omterm{def:qm:finite-difference-methods:theta}{schéma explicite} à $\lambda = 0.489$ (409 pas) donne pour le call à trois mois un prix de 4.3576, contre 4.3576 pour Black--Scholes~; à $\lambda = 0.509$ (393
pas), la fréquence la plus élevée est multipliée par $-1.04$ à chaque pas, et la valeur maximale sur la grille atteint 2\,413.

\section{Stabilité et oscillations}

La stabilité au sens du théorème de Lax ne suffit pas pour les grecques. À la fréquence la plus élevée, le facteur de Crank--Nicolson vaut $(1 - 2\lambda)/(1 + 2\lambda)$~: borné par un, mais proche de $-1$ lorsque
$\lambda$ est grand (la \cref{fig:qm:finite-difference-methods:amp}). Un payoff à angle contient toutes les fréquences, et Crank--Nicolson amortit presque rien les plus hautes, tout en inversant leur signe à chaque
pas. Avec 400 nœuds et 25 pas pour un call à trois mois ($\lambda = 32$), le prix au strike se situe à moins de 0.018 de Black--Scholes, et le gamma oscille entre $-0.81$ et $0.55$ autour d'une valeur réelle de
0.040 (\cref{fig:qm:finite-difference-methods:sawtooth}). Raffiner à la fois l'espace et le temps double $\lambda$ et les dents de scie~: l'erreur maximale sur le gamma vaut 0.46, 0.84, 1.69 et 3.38 avec 200, 400, 800
et 1\,600 nœuds. C'est le lendemain du raffinement.

\begin{omfigure}
\begin{tikzpicture}
\begin{axis}[width=10cm,height=5cm,xlabel={fréquence $\xi\Delta x$},ylabel={facteur d'amplification $g$},
  tick label style={font=\small},label style={font=\small},xmin=0,xmax=3.1416,ymin=-1.15,ymax=1.05,grid=major,grid style={gray!25},
  xtick={0,0.7854,1.5708,2.3562,3.1416},xticklabels={$0$,$\pi/4$,$\pi/2$,$3\pi/4$,$\pi$},
  legend columns=2,legend cell align=left,legend style={font=\footnotesize,at={(0.5,-0.3)},anchor=north,draw=none,fill=none,
  /tikz/every even column/.append style={column sep=8pt}},table/col sep=comma]
\addplot[omThm!50,very thick] table[x=xi,y=explicit049]{figdata/methods/27-finite-difference-methods/amplification.csv};
\addplot[omThm,very thick,dashed] table[x=xi,y=explicit051]{figdata/methods/27-finite-difference-methods/amplification.csv};
\addplot[black,very thick] table[x=xi,y=implicit]{figdata/methods/27-finite-difference-methods/amplification.csv};
\addplot[omDef,very thick] table[x=xi,y=cn]{figdata/methods/27-finite-difference-methods/amplification.csv};
\legend{{explicite, $\lambda = 0.49$},{explicite, $\lambda = 0.51$},{implicite, $\lambda = 32$},{Crank--Nicolson, $\lambda = 32$}}
\end{axis}
\end{tikzpicture}

\omcaption{Facteurs d'amplification de von Neumann des \omterm{def:qm:finite-difference-methods:theta}{theta-schémas} pour $u_\tau = au_{xx}$ en fonction de la fréquence. Le \omterm{def:qm:finite-difference-methods:theta}{schéma explicite} sort de $[-1, 1]$ juste au-delà de $\lambda = \frac12$~; à $\lambda = 32$, le \omterm{def:qm:finite-difference-methods:theta}{schéma implicite}
annule les hautes fréquences et Crank--Nicolson inverse leur signe avec presque aucun amortissement. Données~: le tutoriel du chapitre.}\label{fig:qm:finite-difference-methods:amp}
\end{omfigure}

\begin{omfigure}
\begin{tikzpicture}
\begin{axis}[width=11cm,height=5.4cm,xlabel={prix du sous-jacent $S$},ylabel={gamma},
  tick label style={font=\small},label style={font=\small},xmin=97,xmax=103,ymin=-0.9,ymax=0.65,grid=major,grid style={gray!25},
  legend columns=3,legend cell align=left,legend style={font=\footnotesize,at={(0.5,-0.3)},anchor=north,draw=none,fill=none,
  /tikz/every even column/.append style={column sep=8pt}},table/col sep=comma]
\addplot[omThm!70,thick,mark=*,mark size=1.1pt] table[x=S,y=cn]{figdata/methods/27-finite-difference-methods/sawtooth.csv};
\addplot[black,very thick,dashed] table[x=S,y=exact]{figdata/methods/27-finite-difference-methods/sawtooth.csv};
\addplot[omDef,very thick] table[x=S,y=cn4]{figdata/methods/27-finite-difference-methods/sawtooth.csv};
\legend{Crank--Nicolson,Black--Scholes,{Crank--Nicolson, quatre pas implicites de moitié}}
\end{axis}
\end{tikzpicture}

\omcaption{Gamma d'un call à trois mois à la monnaie obtenu par Crank--Nicolson sur 400 nœuds de log-prix et 25 pas de temps ($\lambda = 32$), avec et sans \omterm{def:qm:finite-difference-methods:rannacher}{démarrage de Rannacher}, au voisinage du strike (les marqueurs sont les nœuds de la grille)~; la dernière courbe se superpose
au gamma de Black--Scholes. Données~: le tutoriel du chapitre.}\label{fig:qm:finite-difference-methods:sawtooth}
\end{omfigure}

\begin{definition}[Démarrage de Rannacher]\label{def:qm:finite-difference-methods:rannacher}
Le \emph{démarrage de Rannacher}\index{démarrage de Rannacher} (Rannacher, 1984) remplace les premiers pas de Crank--Nicolson par des pas d'Euler implicite de moitié, qui amortissent les hautes fréquences de
la donnée initiale non régulière, puis poursuit avec Crank--Nicolson.
\end{definition}

Deux demi-pas implicites suppriment la dent de scie (erreur sur le gamma de $4.7 \times 10^{-4}$ à 400 nœuds), mais laissent une erreur d'ordre un, divisée par deux à chaque raffinement~; quatre demi-pas rétablissent l'ordre deux, l'erreur
diminuant de 4.05 puis 4.07 à chaque division par deux ($1.76 \times 10^{-5}$, $4.35 \times 10^{-6}$, $1.07 \times 10^{-6}$~; la \cref{fig:qm:finite-difference-methods:conv}). Le prix raconte la même histoire, plus discrètement~: Crank--Nicolson pur
ne converge qu'à l'ordre un sur ce problème (erreurs 0.041, 0.018, 0.0092, 0.0046), le démarrage restant à l'ordre deux. Giles et Carter (2006) analysent le mécanisme.

\begin{omfigure}
\begin{tikzpicture}
\begin{loglogaxis}[width=10cm,height=5.4cm,xlabel={nombre de nœuds en espace (pas de temps $= n/16$)},ylabel={erreur max. sur le gamma près du strike},
  tick label style={font=\small},label style={font=\small},xmin=150,xmax=2000,ymin=5e-7,ymax=10,grid=major,grid style={gray!25},
  xtick={200,400,800,1600},xticklabels={200,400,800,1\,600},
  legend columns=2,legend cell align=left,legend style={font=\footnotesize,at={(0.5,-0.3)},anchor=north,draw=none,fill=none,
  /tikz/every even column/.append style={column sep=8pt}},table/col sep=comma]
\addplot[omThm,very thick,mark=*,mark size=1.5pt] table[x=n,y=cn]{figdata/methods/27-finite-difference-methods/convergence.csv};
\addplot[black,very thick,mark=square*,mark size=1.5pt] table[x=n,y=implicit]{figdata/methods/27-finite-difference-methods/convergence.csv};
\addplot[omDef!60,very thick,mark=triangle*,mark size=1.8pt] table[x=n,y=cn2]{figdata/methods/27-finite-difference-methods/convergence.csv};
\addplot[omDef,very thick,mark=diamond*,mark size=1.8pt] table[x=n,y=cn4]{figdata/methods/27-finite-difference-methods/convergence.csv};
\legend{Crank--Nicolson,implicite,{deux demi-pas},{quatre demi-pas}}
\end{loglogaxis}
\end{tikzpicture}

\omcaption{Erreur maximale sur le gamma du call à trois mois pour $90 < S < 110$ lorsque l'espace et le temps sont raffinés ensemble~: Crank--Nicolson pur diverge, le \omterm{def:qm:finite-difference-methods:theta}{schéma implicite} et Crank--Nicolson avec deux
demi-pas implicites convergent à l'ordre un, et avec quatre à l'ordre deux. Données~: le tutoriel du chapitre.}\label{fig:qm:finite-difference-methods:conv}
\end{omfigure}

Le placement du payoff sur la grille compte tout autant. Un digital qui paie 1 au-delà du strike, le strike étant sur un nœud (où le payoff est fixé à $\frac12$), converge à l'ordre un même avec le
démarrage (erreur $9.9 \times 10^{-3}$ sur 200 nœuds, divisée par deux à chaque fois) et ne converge pas du tout sans lui (environ 0.1 à chaque niveau). Avec le strike à mi-chemin entre deux nœuds, ce qui revient à
moyenner le payoff sur la cellule de chaque nœud, l'erreur vaut $6.6 \times 10^{-6}$ sur 200 nœuds et est divisée par 4 à chaque division par deux (\cref{fig:qm:finite-difference-methods:digital}). Pooley, Vetzal et Forsyth (2003)
rapportent le même constat~: moyenner la donnée initiale ou décaler la grille ne suffit pas à lui seul, et ne rétablit la convergence attendue que combiné à un démarrage lissant.

\begin{omfigure}
\begin{tikzpicture}
\begin{loglogaxis}[width=10cm,height=5cm,xlabel={nombre de nœuds en espace},ylabel={erreur sur le prix du digital},
  tick label style={font=\small},label style={font=\small},xmin=150,xmax=2000,ymin=5e-8,ymax=0.03,grid=major,grid style={gray!25},
  xtick={200,400,800,1600},xticklabels={200,400,800,1\,600},
  legend cell align=left,legend style={font=\footnotesize,at={(0.97,0.5)},anchor=east,fill=white,draw=none},table/col sep=comma]
\addplot[omThm,very thick,mark=*,mark size=1.5pt] table[x=n,y=node]{figdata/methods/27-finite-difference-methods/digital.csv};
\addplot[omDef,very thick,mark=square*,mark size=1.5pt] table[x=n,y=midcell]{figdata/methods/27-finite-difference-methods/digital.csv};
\legend{strike sur un nœud,strike entre deux nœuds}
\end{loglogaxis}
\end{tikzpicture}

\omcaption{Erreur sur le prix d'un digital à trois mois au strike, Crank--Nicolson avec quatre demi-pas implicites~: le strike sur un nœud donne l'ordre un, le strike à mi-chemin entre deux nœuds l'ordre deux.
Données~: le tutoriel du chapitre.}\label{fig:qm:finite-difference-methods:digital}
\end{omfigure}

\section{Bords, monotonie et extrapolation}

\begin{definition}[M-matrice, schéma décentré amont]\label{def:qm:finite-difference-methods:m}
Une \emph{M-matrice}\index{M-matrice} est une matrice carrée réelle dont les termes hors diagonale sont négatifs ou nuls, et dont l'inverse existe et est à coefficients positifs ou nuls. Un \emph{schéma décentré amont}\index{schéma décentré amont} discrétise
la dérivée première par la différence décentrée du côté d'où vient la dérive, là où la différence centrée rendrait négatif un coefficient hors diagonale de $L_h$.
\end{definition}

\begin{proposition}[Schémas monotones]\label{prop:qm:finite-difference-methods:m}
Si les termes hors diagonale de $L_h$ sont positifs ou nuls et $c \ge 0$, la matrice du pas implicite $I - \Delta\tau L_h$ est une \omterm{def:qm:finite-difference-methods:m}{M-matrice} pour tout $\Delta\tau$, de sorte que le \omterm{def:qm:finite-difference-methods:theta}{schéma implicite} envoie des données positives ou nulles sur des
solutions positives ou nulles et préserve l'ordre~: un payoff plus élevé donne un prix plus élevé, et aucun nouvel extremum n'apparaît.
\end{proposition}

\begin{proof}
Les termes hors diagonale de $I - \Delta\tau L_h$ sont négatifs ou nuls, et la diagonale de chaque ligne dépasse la somme des valeurs absolues de ses termes hors diagonale (les lignes de $L_h$ somment à $-c \le 0$), de sorte que la matrice
est strictement à diagonale dominante, avec une diagonale positive~: une \omterm{def:qm:finite-difference-methods:m}{M-matrice}, dont l'inverse est à coefficients positifs ou nuls.
\end{proof}

Les différences centrées conservent la structure de signes tant que $|b|\Delta x \le 2a$. Pour un digital à faible volatilité et à taux élevé ($\sigma = 2\%$, $r = 5\%$, $\Delta x = 0.02$), le rapport $b\Delta x/(2a)$ vaut 2.49~; le \omterm{def:qm:finite-difference-methods:theta}{schéma implicite}
avec différences centrées renvoie alors un digital valant 1.030 en un point de la grille, supérieur au paiement actualisé de 0.951, et un prix qui n'est pas monotone en $S$ (variation totale 1.23)~;
le décentrage amont renvoie 0.951 et un profil monotone, au prix d'une précision d'ordre un sur le terme de dérive. Aux extrémités de la grille, le solveur impose soit des valeurs connues (une condition de Dirichlet telle que
$u = S - Ke^{-r\tau}$ pour un call profond dans la monnaie), soit une gamma nulle, avec extrapolation linéaire~; la grille s'étendant sur cinq écarts-types, les deux choix donnent au strike des prix qui diffèrent de $6 \times 10^{-8}$.

\begin{definition}[Extrapolation de Richardson]\label{def:qm:finite-difference-methods:rich}
L'\emph{extrapolation de Richardson}\index{extrapolation de Richardson} (Richardson, 1911) combine deux approximations $u_h$ et $u_{h/2}$ dont l'erreur vaut $Ch^p$ en $u_{h/2} + (u_{h/2} - u_h)/(2^p - 1)$, ce qui
élimine le terme dominant.
\end{definition}

Elle ne fonctionne que si l'ordre est bien celui supposé~: avec le démarrage, les erreurs de prix du call sur 200 et 400 nœuds valent $1.8 \times 10^{-3}$ et $4.6 \times 10^{-4}$, et la valeur extrapolée avec $p =
2$ est à $3.2 \times 10^{-7}$ près. Appliquée à Crank--Nicolson pur, dont l'erreur est ici d'ordre un, elle éliminerait le mauvais terme.

\section{Directions alternées}

\begin{definition}[Méthode des directions alternées implicite]\label{def:qm:finite-difference-methods:adi}
Une \emph{méthode des directions alternées implicite}\index{méthode des directions alternées implicite} (Peaceman et Rachford, 1955~; Douglas et Rachford, 1956) décompose un opérateur multidimensionnel $L = A_0 + A_1 +
A_2$, où $A_1$ et $A_2$ agissent chacun selon une coordonnée et $A_0$ est la partie de dérivée croisée, et traite une direction à la fois de façon implicite. Le schéma de Douglas s'écrit $Y_0 = U^n + \Delta\tau LU^n$, $Y_j =
Y_{j-1} + \theta\Delta\tau A_j(Y_j - U^n)$ pour $j = 1, 2$, $U^{n+1} = Y_2$.
\end{definition}

Chaque étape résout un système tridiagonal par ligne de grille, de sorte qu'un pas coûte $O(n^2)$ sur une grille $n \times n$, au lieu d'une résolution bidimensionnelle. La dérivée croisée est explicite, ce qui rend délicate la stabilité du
schéma~: les conditions sur $\theta$ assurant une stabilité inconditionnelle en présence de termes de dérivée croisée ont été étudiées par Craig et Sneyd (1988) puis par in 't Hout et Welfert (2007), et in 't Hout et
Foulon (2010) comparent les schémas de Douglas, de Craig--Sneyd et de Hundsdorfer--Verwer sur l'équation de Heston (énoncé, non démontré ici)~; la société retient $\theta = \frac12$. Sur une option d'échange $\max(S_1 - S_2, 0)$ avec $\sigma_1 = 30\%$, $\sigma_2 = 20\%$ et une corrélation de 0.5, dont le prix vaut 10.5243 d'après la formule de Margrabe, le solveur de la société
se trompe de 0.045, 0.026 et 0.014 sur des grilles de $40^2$, $80^2$ et $160^2$ nœuds~: un ordre un, parce que la cassure du payoff court le long de la diagonale entre les nœuds, ce que le digital ci-dessus a déjà montré être
le coupable habituel.

\section{Tutoriel~: la dent de scie du gamma}\label{tut:qm:finite-difference-methods:sawtooth}

\begin{tutorial}
\textbf{Objectif.} Pricer un call à trois mois sur une grille de log-prix avec les trois schémas, observer la limite explicite, la dent de scie de Crank--Nicolson et sa suppression. \textbf{État final~:}
les \cref{fig:qm:finite-difference-methods:sawtooth,fig:qm:finite-difference-methods:conv} et les rapports de convergence.
\begin{tutsteps}
\item \textbf{L'opérateur.} Coefficients à trois points sur une grille quelconque, avec décentrage amont optionnel.
\omcode{code/firm/pde/firm_pde.py}{53}{74}{L'opérateur $au_{xx} + bu_x - cu$ sur une grille non uniforme (Python).}
\item \textbf{Le pas de temps et le démarrage}, dans le jumeau C++20 (le jumeau Rust est identique ligne à ligne).
\omcode{code/firm/pde/cpp/firm_pde.hpp}{45}{78}{Pas theta avec conditions aux bords linéaires, et démarrage de Rannacher (C++20).}
\item \textbf{Exécuter} \texttt{explicit\_blowup()}, \texttt{sawtooth()}, \texttt{convergence()}, \texttt{digital\_placement()}, \texttt{stretched()}, \texttt{richardson\_demo()},
\texttt{upwind\_demo()} et \texttt{adi\_exchange()} dans \texttt{qm\_fd.py}~; puis \texttt{fig\_fd.py}.
\end{tutsteps}
\textbf{Pour aller plus loin.} Pricer un call down-and-out dont la barrière est sur un nœud puis entre deux nœuds~; ajouter un \omterm{def:qm:finite-difference-methods:rannacher}{démarrage de Rannacher} à chaque date de surveillance discrète~; essayer $\theta = 0.51$ au lieu du démarrage.
\end{tutorial}

\section{Construction~: le solveur aux différences finies}\label{bld:qm:finite-difference-methods:pde}

\begin{build}
\bfield{Rôle} Équations rétrogrades à un et deux facteurs pour les pricers de la société, avec des grecques pouvant servir à la couverture.
\bfield{Interface} \texttt{uniform\_grid}, \texttt{sinh\_grid}~; \texttt{operator(x, a, b, c, upwind)}~; \texttt{solve\_1d(x, payoff, a, b, c, tau, n\_steps, theta, rannacher, left, right, upwind)}~;
\texttt{thomas}~; \texttt{richardson}~; \texttt{amplification}~; \texttt{interp}~; \texttt{douglas\_adi\_2d}. Jumeaux C++20 et Rust de \texttt{operator\_1d}, \texttt{thomas}, \texttt{theta\_step} et
\texttt{solve\_1d}.
\bfield{Règles} Crank--Nicolson ne démarre jamais sur une donnée non régulière sans demi-pas implicites~; les strikes et les barrières se placent à mi-chemin entre deux nœuds (ou le payoff est moyenné sur la cellule)~; l'ordre se mesure avant
d'extrapoler~; un contrôle de la structure de signes signale une \omterm{def:qm:finite-difference-methods:m}{M-matrice} perdue.
\bfield{Tests d'acceptation} \texttt{code/firm/pde/\{tests,cpp,rust\}}~: Thomas contre une résolution dense~; l'opérateur exact sur les quadratiques~; ordres un et deux mesurés~; la grille étirée bat la
grille uniforme~; condition de Dirichlet~; facteurs d'amplification contre la formule directe~; le décentrage restaure la structure de signes~; ADI contre Margrabe~; C++20 et Rust contre les valeurs Python à $10^{-11}$ près.
\bfield{Pour aller plus loin} Barrières et exercice américain (la méthode de pénalisation pour l'\omterm{def:qm:optimal-stopping-and-impulse-control:fbp}{inéquation variationnelle} du chapitre~10)~; les variantes ADI de Craig--Sneyd et de Hundsdorfer--Verwer~; grilles adaptatives.
\end{build}

\begin{omsources}
\item J.~Crank et P.~Nicolson, \emph{Proceedings of the Cambridge Philosophical Society} 43, 1947~; R.~Courant, K.~Friedrichs et H.~Lewy, \emph{Mathematische Annalen} 100, 1928.
\item P.~D.~Lax et R.~D.~Richtmyer, «~Survey of the stability of linear finite difference equations~», \emph{Communications on Pure and Applied Mathematics} 9, 1956.
\item R.~Rannacher, \emph{Numerische Mathematik} 43, 1984~; M.~B.~Giles et R.~Carter, \emph{Journal of Computational Finance} 9(4), 2006~; D.~M.~Pooley, K.~R.~Vetzal et P.~A.~Forsyth, «~Convergence remedies
for non-smooth payoffs in option pricing~», \emph{Journal of Computational Finance} 6(4), 2003.
\item D.~W.~Peaceman et H.~H.~Rachford, \emph{Journal of SIAM} 3, 1955~; J.~Douglas et H.~H.~Rachford, \emph{Transactions of the AMS} 82, 1956~; I.~J.~D.~Craig et A.~D.~Sneyd, \emph{Computers and
Mathematics with Applications} 16, 1988~; K.~J.~in 't Hout et B.~D.~Welfert, \emph{Applied Numerical Mathematics} 57, 2007~; K.~J.~in 't Hout et S.~Foulon, \emph{International Journal of Numerical Analysis and Modeling} 7, 2010.
\item L.~F.~Richardson, \emph{Philosophical Transactions of the Royal Society A} 210, 1911.
\end{omsources}

\section{Exercices}

\begin{exercise}[$\star$]\label{exo:qm:finite-difference-methods:1}
Avec $\sigma = 20\%$ et $\Delta x = 0.0025$ en log-prix, calculer le plus grand pas de temps explicite stable et le nombre de pas nécessaires à une option à un an.
\end{exercise}

\begin{exercise}[$\star$]\label{exo:qm:finite-difference-methods:2}
Calculer le prix extrapolé, sachant que deux prix obtenus sur les grilles $h$ et $h/2$ valent 4.3558 et 4.3572 et que le schéma est d'ordre deux.
\end{exercise}

\begin{exercise}[$\star$]\label{exo:qm:finite-difference-methods:3}
Montrer que la différence seconde à trois points sur une grille uniforme a pour \omterm{def:qm:finite-difference-methods:consistent}{erreur de troncature} $\frac{h^2}{12}u_{xxxx} + O(h^4)$.
\end{exercise}

\begin{exercise}[$\star\star$]\label{exo:qm:finite-difference-methods:4}
Calculer le facteur de Crank--Nicolson en $\xi = \pi$ pour $\lambda = 32$ et le nombre de pas nécessaires pour amortir ce mode par un facteur 100. Comparer avec un demi-pas implicite.
\end{exercise}

\begin{exercise}[$\star\star$]\label{exo:qm:finite-difference-methods:5}
Pour $au_{xx} + bu_x$, avec des différences centrées sur une grille uniforme, trouver la condition sur $\Delta x$ qui garde positifs ou nuls les deux coefficients hors diagonale.
\end{exercise}

\begin{exercise}[$\star\star$]\label{exo:qm:finite-difference-methods:6}
Expliquer pourquoi une condition de Dirichlet $u = 0$ à la borne inférieure de la grille d'un call ne change presque rien, alors que la même condition à la borne supérieure ruinerait le prix.
\end{exercise}

\begin{exercise}[$\star\star\star$]\label{exo:qm:finite-difference-methods:7}
\emph{Programmation.} Pricer le digital à trois mois, avec le strike sur un nœud puis à mi-chemin entre deux nœuds, par Crank--Nicolson avec quatre demi-pas, sur 200 à 1\,600 nœuds, et donner les rapports de convergence.
\end{exercise}

\begin{exercise}[$\star\star\star$]\label{exo:qm:finite-difference-methods:8}
\emph{Repérer le défaut.} «~Le prix a convergé au centime près quand nous avons divisé la grille par deux, donc les grecques que nous calculons par différences sont justes.~»
\end{exercise}

\section{Problème~: le gamma en dents de scie}

\begin{problem}[{Problème du week-end --- une grecque qui se dégrade avec le raffinement}]\label{pb:qm:finite-difference-methods:1}
Un call à trois mois à la monnaie ($K = 100$, $r = 3\%$, $\sigma = 20\%$) est pricé par Crank--Nicolson sur une grille en log-prix de $\pm5$ écarts-types, avec 400 nœuds et 25 pas de temps.

\medskip\noindent\textbf{Partie I --- Les schémas.}
\begin{enumerate}
\item Quelle est la valeur de $\lambda$, et le \omterm{def:qm:finite-difference-methods:theta}{schéma explicite} serait-il stable~?
\item Où exactement le \omterm{def:qm:finite-difference-methods:theta}{schéma explicite} devient-il instable, et que se passe-t-il juste au-delà~?
\item Quel est le facteur d'amplification de Crank--Nicolson pour la fréquence la plus élevée à ce $\lambda$~?
\item Quelle est la précision du prix au strike, et quelle est celle du gamma~?
\item Que devient l'erreur sur le gamma lorsque l'espace et le temps sont raffinés ensemble~?
\end{enumerate}

\medskip\noindent\textbf{Partie II --- Le remède.}
\begin{enumerate}[resume]
\item Que donnent deux demi-pas implicites, et à quel ordre~?
\item Que donnent quatre demi-pas, et à quel ordre~?
\item À quel ordre le Crank--Nicolson classique converge-t-il sur le prix~?
\item Pourquoi le digital ne converge-t-il qu'au premier ordre lorsque le strike est sur un nœud~?
\item Que donne le déplacement du strike entre deux nœuds~?
\end{enumerate}

\medskip\noindent\textbf{Partie III --- Grilles et extrapolation.}
\begin{enumerate}[resume]
\item Quel gain apporte la grille sinh à nombre de nœuds égal~?
\item Que donne l'\omterm{def:qm:finite-difference-methods:rich}{extrapolation de Richardson} avec le \omterm{def:qm:finite-difference-methods:rannacher}{démarrage de Rannacher}, et pourquoi pas sans~?
\item Quand les différences centrées perdent-elles la propriété de \omterm{def:qm:finite-difference-methods:m}{M-matrice}, et quel en est le coût ici~?
\item Comment se comparent les deux conditions aux limites~?
\item Quelle performance réalise l'ADI de Douglas sur l'option d'échange, et pourquoi seulement au premier ordre~?
\end{enumerate}

\medskip\noindent\textbf{Partie IV --- Jugement critique.}
\begin{enumerate}[resume]
\item Quel test aurait détecté les dents de scie avant que les hedgers ne le fassent~?
\item Pourquoi un \omterm{def:qm:finite-difference-methods:consistent}{schéma stable} ne suffit-il pas pour les grecques~?
\item Quelles règles le pricer de la société devrait-il imposer par défaut~?
\item Énoncer le \emph{résultat nommé}~: le nombre de demi-pas implicites qui supprime l'oscillation et l'ordre qu'il restaure, à partir du rapport d'erreur à chaque division par deux.
\item En une phrase~: que dit le facteur d'amplification à un quant que le prix ne dit pas~?
\end{enumerate}
\end{problem}

\section{Questions d'entretien}

\begin{interviewq}[$\star$ \iqroles{chercheur, développeur}]\label{iq:qm:finite-difference-methods:1}
Quelle est la condition de stabilité du \omterm{def:qm:finite-difference-methods:theta}{schéma explicite} pour l'équation de la chaleur, et pourquoi importe-t-elle pour une grille fine~?
\end{interviewq}

\begin{interviewq}[$\star\star$ \iqroles{chercheur}]\label{iq:qm:finite-difference-methods:2}
Crank--Nicolson est inconditionnellement stable et d'ordre deux. Pourquoi produit-il des grecques oscillantes, et comment le corrigez-vous~?
\end{interviewq}

\begin{interviewq}[$\star\star$ \iqroles{chercheur}]\label{iq:qm:finite-difference-methods:3}
Où placez-vous le strike et la barrière par rapport à la grille, et pourquoi~?
\end{interviewq}

\begin{interviewq}[$\star\star$ \iqroles{développeur}]\label{iq:qm:finite-difference-methods:4}
Comment résolvez-vous efficacement un pas implicite en une et en deux dimensions~?
\end{interviewq}

\begin{interviewq}[$\star\star$ \iqroles{chercheur, risque}]\label{iq:qm:finite-difference-methods:5}
Comment vérifieriez-vous qu'un pricer à différences finies converge à l'ordre qu'il annonce~?
\end{interviewq}

\begin{interviewq}[$\star\star\star$ \iqroles{chercheur}]\label{iq:qm:finite-difference-methods:6}
Votre solveur à deux facteurs donne parfois des prix négatifs pour un digital profondément hors de la monnaie. Que se passe-t-il~?
\end{interviewq}
